\documentclass[11pt]{article}

\usepackage[T1]{fontenc}
\usepackage[utf8]{inputenc}
\usepackage{lmodern}
\usepackage{geometry}
\usepackage{microtype}
\usepackage{xcolor}
\usepackage{setspace}
\usepackage{amsmath,amssymb,amsthm}
\usepackage{enumitem}
\usepackage[round,authoryear]{natbib}
\usepackage{hyperref}
\usepackage{bookmark}

\geometry{margin=1in}
\setstretch{1.08}
\setlength{\parskip}{0.35em}
\setlength{\parindent}{1.25em}
\setlist[itemize]{leftmargin=1.5em, itemsep=0.25em, topsep=0.35em}
\setlist[enumerate]{leftmargin=1.7em, itemsep=0.25em, topsep=0.35em}
\setlength{\emergencystretch}{2em}
\raggedbottom

\definecolor{titleblue}{HTML}{203A5B}
\definecolor{linkblue}{HTML}{1F4E79}

\hypersetup{
  colorlinks=true,
  linkcolor=linkblue,
  citecolor=linkblue,
  urlcolor=linkblue,
  pdftitle={Coherent Possibility and Selective Actualization: A Two-Level Framework for Foundational Explanation},
  pdfauthor={Halvor S. Lande}
}

\newtheoremstyle{foundation}
  {0.85em}
  {0.85em}
  {\itshape}
  {}
  {\color{titleblue}\bfseries}
  {.}
  {0.45em}
  {}

\theoremstyle{foundation}
\newtheorem{theorem}{Theorem}
\newtheorem{proposition}{Proposition}
\newtheorem{definition}{Definition}
\newtheorem{remark}{Remark}

\title{\textcolor{titleblue}{\Large Coherent Possibility and Selective Actualization:}\\[0.35em]\textcolor{titleblue}{\Large A Two-Level Framework for Foundational Explanation}}
\author{Halvor S. Lande}
\date{April 2026}

\begin{document}
\maketitle

\begin{abstract}
A structural or mathematical account can explain admissibility, but it does not yet explain why a determinate world is actual, why actuality is historically articulated, or how becoming proceeds without appeal to a brute selector. This paper argues that the resulting gap is best understood as a problem for foundational explanation. Any adequate foundational account of a nontrivial, dynamically stable, historically articulated world must distinguish two irreducible explanatory roles. The first is constitutive constraint: the conditions under which a candidate structure counts as coherently realizable at all. The second is selective actualization: the conditions under which one admissible continuation rather than another becomes integrated into realized history. I develop the constitutive role in presentation-independent terms through coherent realizability, equivalence-invariance, canonicity, and local satisfiability of coherence burdens. I then propose a localist, regulative characterization of the selective role in terms of commitment, admissibility, local generative affordance, stabilization burden, sufficiency, constructive primeness, and integration. The aim is not to derive a complete ontology or a finished dynamics, but to clarify the explanatory architecture that any adequate foundational explanation must exhibit and to distinguish argued constraints from stronger proposals, unique algorithmic implementations, and deferred formal development.
\end{abstract}

\noindent\textbf{Keywords:} actualization; laws of nature; structural realism; modality; process ontology; philosophy of science

\section{Introduction: the problem of actualization}

Accounts of reality that emphasize structure, mathematics, or modal order promise an explanation of what kinds of world are admissible. They do not yet explain why one admissible continuation is actual, why actuality is historically articulated rather than merely catalogued, or how determinate becoming can proceed without appeal to a brute selector. The gap between admissibility and actuality is the problem addressed in this paper.

The attraction of structural or mathematical metaphysics is nevertheless real. If reality has any intelligible order at all, then it must answer in some way to structure. The hope behind the Mathematical Universe Hypothesis is that the world is not merely \emph{described} by mathematics, but \emph{is} mathematics in a sufficiently robust sense \citep{Tegmark2008}. That ambition belongs to a broader structuralist family in which invariant form does the primary explanatory work \citep{Worrall1989,Ladyman1998,Shapiro1997,French2014,LadymanRoss2007}. The difficulty is that form alone seems to explain, at most, a space of coherent possibility. It does not by itself explain why one possibility is actual, why actuality unfolds through determinate transitions, or why realized history is more than an undifferentiated collection of admissible forms.

The central claim advanced here is therefore not that a final ontology has already been reduced to a definitive pair of laws. It is instead that any adequate foundational explanation of a nontrivial, dynamically stable, historically articulated world must distinguish two irreducible explanatory roles. The first is \emph{constitutive constraint}: it determines what can count as a coherently realizable candidate for existence at all. The second is \emph{selective actualization}: it determines how one admissible continuation rather than another becomes integrated into realized history. The paper argues that these roles are explanatorily distinct even when any later formal theory may articulate them by many principles rather than two single axioms. The proposal is framed at a level prior to familiar disputes over governing, anti-governing, dispositional, and modal accounts of lawhood \citep{Armstrong1983,Dretske1977,Tooley1977,vanFraassen1989,Bird2007,Ellis2001,Mumford2004,Maudlin2007}.

The resulting framework can be stated compactly:

\begin{quote}
\textbf{Constitutive constraint:} the explanatory role that fixes admissible structure.\\
\textbf{Selective actualization:} the explanatory role that governs actual integration into history.
\end{quote}

The second role is \emph{local} in a specific sense that will matter throughout. Locality here is not spatial locality; it is locality to the present continuation cone of realized history. Actualization must proceed by features readable in the realized present and its immediately admissible continuations, rather than by dependence on continuations beyond that field. The proposal is therefore path-dependent without being globally foresighted. As Section 5.1 will make explicit, this locality operates at a level prior to dynamics on a fixed state space, which blocks a natural but mistaken reading of the selective role as a forward-simulating rule.

This paper is a contribution to foundational explanation in philosophy of science, not a complete metaphysical system. It does not introduce a finished evaluator or a cosmology. Its aim is limited: to identify the explanatory architecture required if a determinate world is to be both admissible and actualized, to formulate restrained candidate statements of the two roles, and to distinguish argued constraints from stronger proposals and deferred formal development.

The two laws are offered as regulative characterizations of what any adequate foundational explanation must achieve, rather than as unique algorithmic specifications. This scoping is defended explicitly in subsection 3.3 and, for the selective role in particular, in subsection 5.9.

The immediate target is not every conceivable universe whatsoever, but any world stable enough to preserve determinate identity through change, dynamic enough to exhibit genuine becoming, and rich enough to sustain nontrivial historical articulation. On that target, the paper proceeds as follows. Section 2 situates the proposal among familiar distinctions and identifies the remaining explanatory gap. Section 3 argues that foundational explanation must separate constitutive admissibility from selective actualization. The subsequent sections develop each role in turn, clarify what the framework adds, and conclude by distinguishing what is argued, what is proposed, and what is deferred.

\section{Existing distinctions and the remaining gap}

Several familiar distinctions already mark genuine structure in debates about laws, modality, and world-description. One can distinguish a space of states from a dynamics defined on that space, laws from initial conditions, structural description from realized history, and local evolution from whole-history or global-consistency formulations \citep{Cartwright1983,vanFraassen1989,Maudlin2007,Price1996}. These contrasts are indispensable, but they do not yet resolve the present problem. Even where they explain how a world may be represented, evolved, or globally characterized, they do not by themselves explain why an admissible world is actual or how actualization is articulated from within realized history.

\subsection{State spaces, dynamics, and initial conditions}

The distinction between state space and dynamics captures an important feature of scientific explanation. A state space tells us what kinds of states are being considered; a dynamics tells us how those states evolve \citep{Cartwright1983,Maudlin2007}. But this familiar contrast ordinarily presupposes that the relevant state space has already been specified and that the states within it are intelligible candidates for realization. Constitutive constraint asks a prior question: what makes a candidate structure coherently realizable at all, rather than a merely formal placeholder or a contradictory pseudo-world? For that reason constitutive constraint is not simply a relabeling of state space.

Nor is selective actualization simply a relabeling of dynamics. A dynamical law characteristically maps states to later states, or to trajectories, within an already given framework. By contrast, selective actualization concerns the explanatory role by which one admissible continuation rather than another becomes part of realized history. In a branching setting, or in any setting where multiple admissible continuations remain open, the mere presence of a dynamics does not yet explain actuality. At most it specifies permitted development conditional on an already fixed history. The present proposal therefore cuts across the familiar distinction rather than repeating it. The selective role operates at a level prior to dynamics on a fixed state space: it concerns how admissible continuations become integrated into realized history, thereby contributing to determining what the next state space is, rather than evolving states within an already given one.

The same point holds for the distinction between laws and initial conditions. A law-plus-initial-condition picture may identify a unique trajectory once a starting point is given. But the explanatory burden then shifts to the actuality of that starting point, or to the actuality of the complete boundary data, rather than disappearing \citep{Armstrong1983,Dretske1977,Tooley1977,Bird2007}. If the initial condition is simply stipulated, the transition from admissibility to actuality remains brute. If it is folded into the law, the distinction between constitutive admissibility and selective actualization is still needed to explain what makes the stipulated history both coherent and actual.

\subsection{Structural realism and mathematical description}

Structural realism and related mathematical conceptions of reality rightly emphasize that description should track invariant structure rather than merely superficial presentation \citep{Worrall1989,Ladyman1998,Shapiro1997,French2014,LadymanRoss2007}. That emphasis is closely aligned with one central part of the present proposal. Constitutive constraint is presentation-independent and insists that equivalent descriptions should not be counted as different in being merely because they differ syntactically or coordinatize the same structure differently.

Yet constitutive constraint is not exhausted by structural realism. The present proposal is concerned not only with invariant description, but also with coherent realizability, canonicity, and the local dischargeability of coherence burdens. A candidate world must not only have structure; it must be a structure that can count as realizable rather than permanently schematic, contradictory, or dependent on indefinitely deferred repair. More importantly, structural description alone still leaves open the problem of actuality. Even a fully articulated structural ontology seems, by itself, to describe an admissible space rather than explain why one history within that space is actual.

Tegmark's Mathematical Universe Hypothesis makes the pressure point especially clear. It supplies a constitutive story in terms of mathematical structure, but by treating admissible mathematical structure as already exhaustively real it leaves no selective role to explain why actuality is historically articulated rather than an undifferentiated plurality \citep{Tegmark2008}.

Everettian and branching-multiverse interpretations of quantum mechanics are a distinct class of positions on which admissibility and actuality are not distinguished at the fundamental level. The question of whether such views can reconstruct the target explanandum within a branch, and whether such reconstruction covertly imports a selective role, is a question about physical implementation rather than foundational architecture, and is taken up elsewhere.

\subsection{Primitive actuality and brute selection}

One response to the foregoing difficulty is to stop asking for further explanation and to treat actuality itself as primitive. On that picture, one world, history, or total configuration is simply actual, and no deeper account of the transition from possibility to actuality is needed \citep{Lewis1986,Plantinga1974}. Such a move has the virtue of clarity, and any complete account may eventually have to reach some explanatory terminus. But as a strategy for foundational explanation it leaves too much of the present problem untouched.

If actuality is treated as primitive in that strong sense, then the selection of one admissible continuation rather than another is no longer illuminated but merely named. The present framework does not deny that actuality is basic in the sense of being irreducible to mere description. It argues, rather, that the explanatory role of actuality can still be constrained. Selective actualization is proposed precisely as an account of what such constraint must look like if actuality is not to enter the story as an unconstrained brute selector. In that respect the proposal is closer to recent grounding-based attempts to articulate dependence relations without collapsing them into simple reduction \citep{Fine2012,Schaffer2009,Rosen2010}. The role must remain tied to admissibility, path-dependent history, local generative affordance, and integration into what follows, and it must do so at the explanatory level appropriate to a role that fixes, rather than presupposes, the next stage of history.

\subsection{Whole-history and global-consistency views}

Whole-history and global-consistency views identify another genuine insight. The actual world may not be understandable one step at a time if large-scale consistency conditions, variational principles, or complete-history constraints are constitutively relevant \citep{Price1996,WheelerFeynman1945,WheelerFeynman1949,Cartwright1999}. The present proposal does not deny that possibility. A global constraint may well help determine what counts as an admissible history, and in that sense such views may contribute to the constitutive side of foundational explanation.

What they do not yet provide, at least not by themselves, is an account of local articulation. If actuality had to proceed by surveying complete histories as complete histories before any local becoming could occur, then the explanation of becoming would become either teleological in formulation or opaque in operation. The localizability requirement defended later in the paper is meant to block that result without caricaturing holistic alternatives. The claim is not that global constraints are irrelevant; it is that any adequate account of an actual, historically articulated world must also explain how actuality is readable from within the present continuation cone rather than only from a finished God's-eye view of total history.

\subsection{The governing/anti-governing debate and the level of the present proposal}

The distinctions canvassed so far are distinctions within philosophy of science about how laws should be understood once one already has a world in which laws can be stated. A further debate concerns the ontological status of laws themselves. Governing accounts hold that laws are entities or relations distinct from the particulars whose behavior they constrain, and that laws genuinely govern such behavior by standing in a primitive necessitation relation to it \citep{Armstrong1983,Dretske1977,Tooley1977,Maudlin2007}. Anti-governing accounts, most influentially the Humean best-system tradition, hold that laws are nothing over and above the most efficient systematization of the particular facts that make up the history of the world \citep{Lewis1986,Loewer1996,Hall2015}. Dispositional essentialism, by contrast, grounds laws in the natures or essences of properties and objects, so that the regularities laws describe express what it is to be the kinds of things the world contains \citep{Bird2007,Ellis2001,Mumford2004}.

The two-role framework developed in this paper is deliberately prior to this debate. It is not offered as a fourth position alongside governing, anti-governing, and dispositional accounts. It is a proposal about the explanatory architecture that any such position must exhibit if it is to account for a determinate, historically articulated world.

The reason can be stated directly. Each of the three positions is a proposal about what the relation is between laws and particulars within realized history. The Humean says the relation is one of summarization. The governing theorist says it is one of necessitation. The dispositional essentialist says it is one of essential grounding. But all three positions presuppose, and do not by themselves explain, two prior questions. What makes something a candidate for being a particular in the first place, what distinguishes a coherent, realizable structure from a pseudo-structure that could not be a particular at all? And what determines which of the admissible candidates become integrated into realized history rather than another, what distinguishes the particular facts that constitute the world's history from the ones that remain merely possible?

These are the constitutive and selective questions respectively. Each of the three positions takes for granted that some answer to both is available. The Humean presupposes a history of particular facts to be systematized, which presupposes both that such facts are coherent candidates for being and that they have become actual rather than merely possible. The governing theorist presupposes that there are particulars for laws to govern, which likewise presupposes answers to both prior questions. The dispositional essentialist presupposes that there are property-instances with essences, which again presupposes both coherent admissibility and selective actualization of those instances.

This clarifies the relation to Maudlin in particular. Maudlin is a near neighbor because he rightly insists that laws must do more than summarize already-given facts and must help account for determinate history \citep{Maudlin2007}. The present proposal departs from his view at two points. It separates constitutive constraint from selective actualization, whereas Maudlin's primitivism does not sharply divide those explanatory burdens. And it locates selective actualization at a level prior to dynamics on a fixed state space, rather than treating primitive law as a governing relation internal to an already given state space with a fundamental temporal direction.

It follows that the two-role framework neither competes with nor depends on any particular answer to the governing debate. It operates at a level prior to the debate by articulating what the debate takes for granted. A Humean, a governing theorist, and a dispositional essentialist can all accept the two-role architecture while disagreeing about what laws are within realized history, because the two-role architecture concerns the conditions of there being realized history at all rather than the internal structure of the laws that describe it.

This positioning has a further consequence worth making explicit. It is sometimes said against the Humean tradition that it cannot explain why one particular history is actual rather than another, because the best-system account only systematizes the history that in fact obtains and does not explain the obtaining \citep{Maudlin2007}. The two-role framework clarifies what this complaint is and is not. If the complaint is that the Humean cannot locate the explanation of actualization within the theory of laws, then the two-role framework agrees but generalizes the point: no theory of laws can locate that explanation within itself, because the selective role is a distinct explanatory role operating at a prior level. The governing theorist is not in a better position than the Humean on this question, despite appearances. Primitive necessitation between universals and particulars presupposes that there are particulars to be necessitated, which is already a fact about selective actualization rather than about governing.

The same reframing applies to the debate over whether laws are metaphysically prior to particulars or vice versa. Both sides presuppose a distinction between what is admissible and what is actual, and both sides presuppose that this distinction is articulated by some structural architecture. The two-role framework makes that architecture explicit and proposes that any adequate metaphysics of laws, Humean, governing, dispositional, or otherwise, must exhibit it.

\subsection{Scope and presentation-independence}

The remaining gap can now be stated more precisely. The problem is not just how to distinguish possibilities from actualities in the abstract. It is how to explain the architecture by which coherent admissibility and realized actuality belong together in one determinate world. Constitutive constraint concerns what can count as a coherently realizable candidate for existence in the first place. Selective actualization concerns the integration of one admissible continuation rather than another into realized history. Those questions arise at the level of foundational explanation rather than at the level of already-given physical models.

The present argument therefore concerns foundational explanation rather than any privileged implementation. It is not the claim that reality is literally ``written'' in any one human formalism, still less that the world is a proof assistant or that its source code can be identified with a particular syntactic calculus. Human logical systems are at best presentations of deeper invariant content. They matter here only insofar as they provide partial insight into what constitutive admissibility and selective actualization would have to preserve. Constructive type theory is useful because it treats existence as witness-bearing and because it makes the passage from possibility to actuality formally visible in the judgmental form $\Gamma \vdash a : A$ \citep{MartinLof1984}. Homotopy type theory is useful because it exhibits a presentation in which identity carries internal structure and equivalence can be elevated to ontological significance \citep{UF2013}. But the paper does not assume that reality is identical with either presentation. It seeks the invariant content beneath them.

Talk of a foundational law in what follows is shorthand for a role in foundational explanation, not yet a contingent regularity of a particular physical cosmos. The paper does not introduce a specific evaluator, local-actualization rule, or cosmological mechanism. It argues only that foundational explanation must separate two roles and that these roles must have a certain structural shape. It also does not treat arithmetic as a third primitive; stagewise language is used only as a weak convenience for speaking about cumulative history.

\section{A Two-Level Framework for Foundational Explanation}

\subsection{The insufficiency of constitutive order alone}

Suppose there were only a constitutive order: a realm of coherent forms, structures, or possibilities. Then one of three things must happen. Either nothing becomes actual, in which case no world exists. Or every coherent possibility is equally actual, in which case actuality is indistinguishable from undifferentiated possibility and there is no determinate history. Or one possibility is actual rather than another for no reason internal to the ontology, in which case actuality is brute.

None of these options is satisfactory. The first yields no world. The second yields no articulated world. The third abandons explanation precisely where ontology should begin. A constitutive order alone is therefore incomplete.

\subsection{The insufficiency of actualization alone}

Now suppose there were only an order of actualization. An actualizing principle presupposes something to be actualized. If there is no prior distinction between coherent and incoherent, admissible and inadmissible, meaningful and meaningless, then the very idea of actualization is empty. An ``actualization'' that can equally well realize contradiction, incoherence, or indeterminacy is no law of reality at all.

Hence actualization cannot be primitive in isolation. It presupposes a prior order of coherent possibility.

These paired insufficiencies support the paper's core thesis. Any adequate foundational explanation of a determinate and historically articulated world must distinguish at least two irreducible explanatory roles: one that constrains what can coherently count as admissible structure, and one that governs selective actualization from among admissible continuations. The claim is architectural rather than axiomatic. It concerns the explanatory work that must be done, not the final number of principles in a completed formal theory.

\subsection{Why putative third roles do not add a new explanatory level}

The claim that, for the target explanandum fixed in Section 1, foundational explanation requires exactly two irreducible roles is strong, and it invites a natural objection. Could there not be a third primitive role, a role of modal structure, of individuation, of normativity, or of something else, that operates at the same foundational level as constitutive constraint and selective actualization without reducing to either? If such a role is possible, the architectural thesis of Sections 3.1 and 3.2 is incomplete. This subsection argues that the two-role architecture is exhaustive for that indexed explanandum, and clarifies what that indexed exhaustiveness claim does and does not assert.

The argument proceeds in three stages. First, it characterizes what foundational explanation is trying to do and shows that the explanandum decomposes into exactly two irreducible questions. Second, it identifies three categories into which putative third roles fall on examination: structural consequences of how the two roles operate, derived features of cumulative actualization, and scope questions that lie outside foundational inquiry rather than expanding its internal architecture. Third, it tests this classification against the most pressing candidates, above all modal structure and individuation, while also indicating how temporal directionality, probabilistic refinement, and observer-style proposals are to be located.

\subsubsection{The structure of the explanandum}

Foundational explanation, as understood here, asks what must be the case for there to be a determinate, historically articulated world at all. This question is not a question about the content of physics, about the structure of any particular domain, or about which features of reality are contingent. It is a question about what any world of the relevant kind must exhibit in order to be such a world. The question is therefore prior to the specification of any particular ontology and answers to no empirical inquiry.

So characterized, the explanandum decomposes into two irreducible sub-questions.

The first concerns admissibility. What makes something a candidate for being at all? What distinguishes a coherent, realizable structure from a pseudo-structure that only seems to specify a possible way for things to be? This is the constitutive question, and Section 3.1 argued that any adequate account must answer it.

The second concerns actualization. Among the candidates that pass the first filter, what determines which are integrated into realized history? What distinguishes one admissible continuation's becoming actual from another's? This is the selective question, and Section 3.2 argued that any adequate account must answer it as well.

The indexed exhaustiveness thesis is that these two questions exhaust the foundational explanandum as that explanandum has been delimited here. There is no third foundational question a third primitive role could answer to within that target. This is not a claim that metaphysics contains only two topics; metaphysics contains many topics. It is a claim that at the level at which constitutive and selective roles operate, there are exactly two irreducible questions, and any further question either reduces to a sub-question of these, emerges as a structural feature of how they operate, or lies outside the scope of foundational explanation altogether.

This thesis is indexed to a particular conception of what foundational explanation is. A philosopher who takes normativity, teleology, intentionality, or value to be foundational in the sense relevant here would reject the thesis. The present argument does not attempt to defeat that philosopher on their own terms. It claims only that, on the conception of foundational explanation operative in this paper, a conception concerned with the conditions of there being a determinate, historically articulated world rather than with the normative, axiological, or intentional structure of such a world, the two-question decomposition is exhaustive.

This indexing is not an evasion. All foundational philosophy operates from some characterization of what foundational inquiry is for, and Kant's distinction between constitutive and regulative principles, as reconstructed in Friedman's work on the relativized a priori, is the paradigmatic acknowledgment that foundational claims are indexed to a prior characterization of the domain to be grounded \citep{Friedman2001}. The thesis developed here is indexed to a characterization it shares with much of the structural realist tradition, which treats the question of what structure a determinate world must exhibit as the proper locus of foundational inquiry \citep{LadymanRoss2007,French2014}.

This indexed formulation also helps locate several familiar candidates for a third role. A fundamental temporal direction of actualization, if posited, does not answer a third question over and above admissibility and actualization. It specifies an asymmetry in how the selective role operates: realized history is preserved and extended rather than erased or actualized backwards. Likewise a probabilistic measure, amplitude rule, or tie-breaking distribution over equally eligible continuations does not answer a third foundational question. It is a more articulated answer to the selective question, because it refines how one admissible continuation rather than another becomes actual once the foundational non-dominance conditions are fixed. Appeals to observers or consciousness divide in the same way. If consciousness helps determine what counts as an admissible candidate, it belongs on the constitutive side. If it helps determine which admissible continuation is realized, it belongs on the selective side. Only if one insists that consciousness, normativity, or value is a third irreducible kind of foundational explanation does the present thesis cease to apply, and that is exactly the indexed disagreement just noted.

\subsubsection{How putative third roles are to be classified}

Given the two-question characterization of the explanandum, any candidate for a third primitive role must either answer one of the two questions more specifically, redescribe the results of their interaction, or stand outside the framework altogether. This yields a three-part classification of candidate third roles.

The first category consists of roles that, on examination, are \emph{structural consequences} of how constitutive constraint and selective actualization operate rather than additional primitive roles alongside them. Such roles describe real features of the two-role architecture, but they are not themselves foundational primitives. Modal structure, examined below, falls into this category.

The second category consists of roles that are \emph{derived features of cumulative actualization}, roles that describe aspects of realized history as it accumulates, and that are explained by the path-dependent operation of the selective role together with the constitutive constraints on what can be realized. Such roles are real features of historically articulated worlds, but they are not primitive because they are explained by the cumulative character of selective actualization already captured in Section 5.3 under the notion of commitment. Individuation, examined below, falls into this category.

The third category consists of \emph{scope questions}, questions that appear to demand a foundational role but that actually lie outside the scope of foundational explanation as characterized in subsubsection 3.3.1. These are questions whose answers would presuppose the very roles they are trying to explain. Asking what grounds the constitutive role would require an explanation operating prior to the conditions of coherent realizability, which is incoherent because any explanation presupposes such conditions. Asking what grounds the selective role would require an explanation operating prior to the conditions of actualization, which is likewise self-undermining because any explanatory move is itself a move within the space of actualizable claims. Questions of this kind are not left unanswered because the framework is incomplete; they are recognized as standing outside the framework in a principled sense, for reasons that are themselves articulable within foundational philosophy.

This third category deserves direct attention because it is often the source of the impression that the two-role architecture is incomplete. The impression arises when one expects foundational explanation to ground not only the structure of any determinate world but also the possibility of there being such structure at all. That further demand is intelligible, but it is not a demand that can be met by adding a third role. A third role would still be one constituent within the architecture and would therefore still presuppose the very framework whose instantiation is being asked about. The residual demand is instead for an account of the instantiation of the architecture rather than for another within-framework primitive. Later Wittgenstein and contemporary structural realism make adjacent points about boundary propositions and stopping-points of inquiry, but the present claim does not stand or fall with those analogies \citep{Wittgenstein1969,LadymanRoss2007,French2014}.

The two-role thesis therefore does not claim that the framework explains everything that might be called a foundational matter. It claims that at the level of foundational explanation, properly scoped, exactly two irreducible roles are needed, and that apparent demands for a third role resolve into one of the three categories above.

\subsubsection{The case of modal structure}

Modal structure is the most natural candidate for a third primitive role. A critic might say: constitutive constraint fixes what is admissible, and selective actualization governs what becomes actual, but neither tells us how admissible alternatives \emph{relate to one another as alternatives}. That relation, the modal texture of the space of admissibility, appears to be a primitive structural feature, neither reducible to admissibility itself nor to the selective operations that integrate some admissible candidates into history.

This is a substantive challenge, but modal structure does not in fact require a third primitive role. To see why, consider what it is for two admissible candidates to be alternatives to one another. The relation has two components. The first is that both candidates are admissible, that they pass the constitutive filter. The second is that their joint realization is excluded by how actualization operates, that integrating one into history excludes the integration of the other at the same stage. The first component is constitutive. The second is selective. The modal relation of being an alternative to is therefore not a further primitive fact over and above these two components; it is exactly the conjunction of a constitutive fact and a selective fact.

This is the first category of the classification: modal structure is a structural consequence of how the two roles operate, not an independent primitive. The point generalizes. Richer modal relations, relations of comparative accessibility, of relative necessity, of counterfactual dependence, are similarly analyzable as structural features of the admissibility space and the selective operations that determine integration. A putative primitive modal role would have to attach to some modal fact not captured by constitutive plus selective resources, and the burden on the proponent of such a role is to specify what that fact is. Lewis-style modal realism, Finean grounding-theoretic approaches, and dispositional essentialist treatments of modality have each attempted to identify such a fact, and each has faced the persistent difficulty that the modal facts they propose turn out either to be analyzable in terms already available within a constitutive-plus-selective architecture or to presuppose such an architecture in order to be stateable \citep{Lewis1986,Fine2012,Bird2007}.

The thesis is therefore not that modal structure is unimportant or that modal questions are misguided. It is that the work done by modal structure at the foundational level is work that constitutive and selective roles, operating together, already do. Adding a third primitive modal role would be redundant rather than explanatory.

\subsubsection{The case of individuation}

A harder candidate is individuation. The Selective Law as developed in Section 5 operates on a particular realized history $H$ and its admissible continuation cone $A(H)$. But what makes \emph{this} continuation cone rather than another the relevant one? Why is $H$ the realized history to which the selective role is applied, rather than some other admissible history? The question appears to call for a third role, a role of individuation that selects which realized history is the subject of further actualization in the first place.

This candidate repays careful examination. The individuation question presupposes that there is a ``this'' realized history that has been distinguished from alternatives, and it asks what distinguishes it. The answer available within the two-role framework is that $H$ is distinguished as the cumulative product of prior applications of the selective role to prior realized histories, tracing back through the path-dependent accumulation of actualization. At each prior stage, one admissible continuation was integrated rather than another, by the operation of the selective role. The present realized history is nothing over and above this cumulative path. No third role is required to individuate it, because its individuation is exhausted by the cumulative, path-dependent character of selective actualization itself, what Section 5.3 identifies as commitment.

This is the second category of the classification: individuation is a derived feature of cumulative actualization, not an independent primitive. The relevant continuation cone is this one because the selective role has operated in the particular path-dependent way it has, and commitment, the preservation of realized history as a determinate context for further selection, secures the individuation of $H$ without requiring additional primitive structure.

But the examination is not yet complete, because a sharper version of the individuation worry remains. Even granting that each realized history individuates later continuation cones via commitment, what individuates the \emph{first} realized history, the starting point of the whole path-dependent process? If individuation reduces to commitment, and commitment requires something to commit to, then the origin of the first realized state appears to require something beyond what the two roles can provide.

This is where the initiation question has to be split rather than merely named. One version asks what could count as an admissible initial stage. That is a constitutive question. Another asks, if more than one initial stage is admissible, what could order their realization without global survey or brute stipulation. That is a selective question, simply posed at a boundary case in which the inherited history is minimal. In that limited but important sense the framework does answer the initiation question: any initial actuality must lie within the constitutively admissible space, and any nontrivial choice among admissible initial candidates must be governed by the same kind of selective discipline that governs later stages. There is no third task of ``getting actualization started'' over and above fixing admissible initial candidates and ordering their realization.

What remains is the genuine boundary question why there is any realized history at all rather than none, or why one admissible initial candidate rather than another is realized when the implementation-level refinement of the selective role is left unspecified. That residue should not be denied. But it is not evidence for a third primitive role. The first part asks for the existence of actuality rather than for the internal architecture of actualization. The second asks for a more determinate articulation of the selective role than the foundations paper undertakes. Both outrun the present target explanandum, though in different directions.

The honest treatment of initiation is therefore neither to pretend the framework answers every boundary question nor to infer a new role from the leftover residue. It is to mark a limiting case. The framework constrains initial actuality from within, but it does not purport to explain why the architecture is instantiated at all. That residual question is principled because answering it would require a theory of the instantiation of the architecture, not a third constituent within the architecture itself. The later Wittgenstein is at most corroborative here rather than load-bearing: not every legitimate boundary condition becomes a further within-framework primitive \citep{Wittgenstein1969}.

\subsubsection{The shape of the exhaustiveness thesis}

Taken together, the preceding examinations support a more precise statement of the exhaustiveness thesis than subsection 3.3 as originally formulated offered. The thesis is this. At the level of foundational explanation characterized in subsubsection 3.3.1, and for the target kind of world fixed in Section 1, the architecture of any adequate account must exhibit exactly two irreducible roles, one answering to each of the two sub-questions into which the foundational explanandum decomposes. Apparent demands for a third role admit of a principled three-part classification: some candidate roles are structural consequences of how the two roles operate, some are derived features of cumulative actualization, and some are scope questions lying outside foundational explanation rather than expanding its internal architecture. The two cases examined, modal structure and individuation, exemplify the first and second categories respectively, and the discussion of initiation illuminates the third.

The classification is not a proof that no future candidate third role could evade it. It is a structural proposal and a stress-tested schema: that any candidate third role, when properly examined, will fall into one of the three categories, and that no candidate has yet been offered that succeeds in occupying a genuine fourth position. The thesis is defended by the structural argument of subsections 3.3.1 and 3.3.2 rather than by exhaustive enumeration, and its vulnerability is to the production of a putative foundational role that cannot be classified into any of the three categories. The burden on the proponent of such a role is to specify what foundational question the role answers to and to show that the question is neither a sub-question of admissibility or actualization, nor a structural consequence of their operation, nor a scope question about the framework itself.

Until such a role is articulated, the two-role architecture is the best-supported architecture for the indexed explanandum identified here. That is the sense in which the thesis is exhaustive. It is an exhaustiveness claim about the explanatory architecture of a delimited target, not an unconditional guarantee that no richer metaphysical outlook could adopt a wider foundational basis.

\begin{remark}[Explanatory level]
Both roles identified here operate at the level of foundational explanation rather than at the level of dynamics within an already specified state space. The constitutive role fixes what counts as an admissible candidate; the selective role governs which admissible continuation becomes integrated into realized history. Neither is a rule for evolving states within a pre-existing arena. This is not a terminological point. It affects how the two roles should be assessed: objections that presuppose a state-space-plus-dynamics framing may succeed against proposals at that level without touching the present thesis. Section 5.1 will return to this point in connection with the locality of the selective role.
\end{remark}

\section{Constitutive Constraint}

\subsection{Why possibility must be coherent}

The minimal requirement on a possible world is not merely that it be thinkable, but that it be \emph{coherently realizable}. An entity or process cannot count as genuinely possible if its very being depends on unresolved contradiction. A square circle does not fail because it is unusual; it fails because the conditions required to specify it cancel one another.

Yet bare non-contradiction is too weak. A world is not constituted by a list of propositions that happen not to clash. It is constituted by structures that must compose, persist, and interact. Constructive approaches are helpful here because they refuse to identify existence with mere verbal consistency and instead require something closer to witness-bearing realizability \citep{MartinLof1984,UF2013}. Thus possibility must be understood as \emph{coherent realizability}: the mutual satisfiability of the structure's internal and external conditions of being.

\begin{definition}[Coherent realizability]
A candidate structure is coherently realizable when its constituents, relations, and operations can be jointly determined without unresolved contradiction or indefinite suspension of compatibility conditions.
\end{definition}

\subsection{Identity, equivalence, and invariance}

A second constitutive requirement concerns identity. If reality is structured, then sameness cannot be treated as a merely flat predicate detached from the internal relations that make a thing what it is. The identity of a structure must be rich enough to support the preservation and transport of its organization across admissible variation.

This does not yet force any one formal doctrine of identity, but it does force a principle of \emph{equivalence-invariance}. If two presentations differ only in ways that make no invariant structural difference, then the difference cannot be ontologically fundamental. Reality cannot count as different what differs only by representational dressing. In that respect the present argument aligns with structural realist and mathematical structuralist insistence that ontology should answer to invariant organization rather than accidental encoding \citep{Worrall1989,Ladyman1998,French2014,Shapiro1997}.

\begin{definition}[Equivalence-invariance]
A constitutive order is equivalence-invariant when representationally distinct but structurally equivalent candidates are not counted as different in being.
\end{definition}

This principle is important because without it ontology becomes hostage to presentation. The same world would be multiply overcounted merely by redescribing it. A genuine constitutive law must cut deeper than syntax.

\subsection{Canonicity and executability}

A third requirement is canonicity. A possible world cannot be built from permanently stuck pseudo-objects whose closed realization never becomes determinate. If a construction is genuinely closed, reality must determine what it is. Otherwise the candidate remains an unresolved schematic placeholder rather than an executable world-constituent. This is one reason constructive formalisms remain philosophically instructive even when the present paper does not identify reality with any one of them \citep{MartinLof1984,UF2013}.

\begin{definition}[Canonicity]
A constitutive order satisfies canonicity when closed constructions admit determinate realization rather than remaining indefinitely non-canonical.
\end{definition}

The language of executability is not meant computationally in any narrow machine sense. The point is ontological: a world must be capable of running as a world. A purely schematic form that never resolves into determinate being is not enough.

\subsection{Local satisfiability of coherence burdens}

A fourth requirement concerns scale. Any genuine structure carries compatibility conditions. If these coherence burdens can grow without local discharge, then no stable world is possible. There would be only an ever-accumulating debt of unresolved compatibility, and reality would collapse into non-executability.

Hence a constitutive law must require that coherence burdens be \emph{locally satisfiable}. New structure may introduce new conditions, but those conditions must be dischargeable within the world rather than indefinitely deferred.

\begin{definition}[Local satisfiability]
A constitutive order satisfies local satisfiability when the compatibility burdens introduced by new structure are locally dischargeable rather than requiring indefinitely unresolved global repair.
\end{definition}

\subsection{Canonical formulation}

We may now state the constitutive law in its strongest compact form.

\begin{theorem}[Constitutive Law]
A universe of the target kind can exist only as a coherent, equivalence-invariant, canonically executable structure whose identity and compatibility conditions are locally satisfiable.
\end{theorem}

This statement is intentionally presentation-independent. It does not identify the law with any one proof theory or any one formal semantics. Rather, it names the invariant content that any adequate presentation must preserve.

\subsection{Why weaker alternatives fail}

The constitutive law has the shape above because weaker candidates fail.

Bare consistency is too weak. One can imagine a body of non-contradictory descriptions that nevertheless never becomes an executable world because its constructions remain non-canonical or its compatibilities remain unresolved.

Flat identity is too weak. If sameness is treated as a purely extensional yes-or-no fact detached from structural transport, ontology becomes blind to the internal articulation of identity and loses the ability to distinguish genuine sameness from accidental coding.

Representation-sensitive ontology is too weak. If structurally equivalent presentations count as different worlds, being becomes hostage to syntax.

Unbounded coherence debt is too weak. If new structure may indefinitely accumulate unresolved compatibility burdens, then no stable world can be realized.

Taken together, these failures force the stronger constitutive statement above.

The constitutive law should also be read regulatively rather than algorithmically. It characterizes what any adequate foundational constitution must achieve, coherent realizability, equivalence-invariance, canonical executability, and local satisfiability, without claiming that a single formalism or evaluator uniquely implements those requirements. In that respect it stands in deliberate parallel with the selective law developed in Section 5.

\section{Selective Actualization}

Constitutive constraint yields a realm of coherent possibility. It does not yet yield reality as realized history. The task of selective actualization is therefore not to replace the constitutive order, but to govern how actuality proceeds from within it.

Let $H$ denote realized history at some stage. Let $A(H)$ denote the admissible continuation cone at $H$: the locally accessible continuations compatible with constitutive constraint and with the commitments already carried by $H$. Selective actualization is a rule on $A(H)$. The question is: what structural features must such a rule possess if it is to generate a nontrivial, dynamically stable, historically articulated world without hidden global foresight?

\subsection{Why actuality must be local}

A determinate universe cannot require an external oracle that surveys indefinitely remote futures before each act of becoming. If actuality depended on such a survey, the law of becoming would either be teleological in formulation or opaque in operation. Time-symmetric and whole-history approaches sharpen the temptation to such a picture, and also the cost of yielding to it \citep{Price1996,WheelerFeynman1945,WheelerFeynman1949}. If Price is right that the deepest physical description is time-symmetric, that sharpens rather than weakens the present point: the articulation of realized history from within the present continuation cone is then explanatory work that the time-symmetric base does not itself perform. Any viable Selective Law must therefore be \emph{localizable}: it must order candidate continuations by features available in the realized present and its immediately accessible continuation cone.

Two clarifications are needed before this requirement can be stated precisely, because the concept of locality at stake is easily confused with nearby but distinct notions.

First, locality here is not spatial. A continuation cone is a cone of admissible extensions of realized history, not a region of space or spacetime. The requirement constrains what the selective role may presuppose, not where its effects may be felt.

Second, and more importantly for what follows, locality here does not operate at the level of dynamics on a fixed state space. This distinction matters because the most natural way to raise an objection against localism, namely, the worry that the selective law covertly looks ahead into the future, presupposes that the selective role is a dynamics of that kind. It is not. A dynamics in the familiar sense presupposes a state space already given, within which forward evolution is defined. Selective actualization, as articulated here, operates at a level prior to such dynamics: it concerns the integration of admissible candidates into realized history, and thereby contributes to determining what the next realized history is. The distinction between constitutive and selective roles developed in Section 3 is, among other things, a distinction of explanatory level. Selective actualization is not a forward-simulating rule running on a prefabricated stage. It is one of the conditions under which a determinate next stage becomes available at all.

The remaining question, taken up in subsection 5.2, is which structural features of the present continuation cone the Localizability Requirement may permit the selective law to read.

\subsection{The Localizability Requirement}

\begin{definition}[Localizability Requirement]
Any viable Selective Law must order admissible continuations by properties encoded in the realized history it inherits, in the candidate continuation itself, and in the immediate structural relations between the candidate and the admissible field opened by its hypothetical integration.
\end{definition}

Three features of this formulation deserve comment.

First, the properties that ground comparative eligibility must be readable from structure already determined at the present stage. Realized history is already determinate. A candidate continuation is, by hypothesis, a specific admissible extension under consideration. The relations between a candidate and its immediate continuation field are structural features of the candidate's position in the admissible order, not forecasts of what will later happen.

Second, admissible comparisons may include \emph{structural counterfactual} properties of the candidate. What a candidate's integration would make available, as a matter of admissibility structure, is a feature of the candidate's place in the admissible order at the present stage. It is not a prediction about which admissible continuations will \emph{in fact} be integrated at subsequent stages. The distinction is the same distinction already standard in discussions of modal structure: to say that a move opens certain further admissible moves is to describe its position in the structure of admissibility, not to predict a trajectory through it.

Third, the Localizability Requirement forbids dependence on properties that are not readable from the present continuation cone, in particular, dependence on how far-downstream admissible developments would actually unfold, or on the content of continuations that do not figure in the immediate admissible field. This is what blocks teleology and computational opacity.

Selective actualization is therefore path-dependent without being globally foresighted. It reads what is already structured into the present admissible field. It does not simulate the future on a pre-existing stage, because at the level at which it operates there is no pre-existing stage on which such simulation could be defined.

\subsubsection{Addressing the look-ahead objection}

A natural objection to the framework deserves direct response. The Selective Law, as stated below in its canonical formulation, appears to rank continuations partly by their \emph{intrinsic generative affordance}, a measure of how much further admissible continuation space a candidate immediately opens. Does this smuggle in a covert look-ahead, so that localism is preserved in letter but not in substance?

The objection has force if selective actualization is read as a dynamics on a fixed state space. On that reading, to assess what a candidate ``opens'' is to forecast which future states will be reachable from it, which requires simulating the dynamics forward. That is a global operation, and the Localizability Requirement would indeed fail to block it.

But the reading misidentifies the explanatory level at which the selective role operates. The distinction drawn in the previous subsection is relevant here. Intrinsic generative affordance is not a forecast of which admissible continuations will subsequently be integrated. It is a structural property of the candidate's position in the admissible order at the present stage: the immediate admissible field its integration would open, read from features already present in the candidate's specification and its relation to realized history. Whether an admissible opening will in fact be taken at some later stage is a separate question, settled by further applications of the Selective Law to further realized histories. It is no part of what the current application needs to determine.

An analogy from a different domain may help, though it is only an analogy. When a constructive proof introduces a lemma, one can ask what further inferences that lemma makes immediately available. The answer is a structural property of the lemma's position in the inferential order, readable from the lemma itself together with the prior context, and does not require predicting which further inferences will in fact be drawn. The distinction between what a move structurally opens and what will trajectorially follow from it is not peculiar to the present framework; it is already operative wherever structural and dynamical notions are distinguished.

Two further points secure the position.

First, the Selective Law compares candidates at a common stage. Comparison requires reading each admissible candidate and the immediate structural relations among them. This is a contemporaneous operation on the present admissible field, not a forward projection. The Localizability Requirement permits such contemporaneous comparison; what it forbids is dependence on continuations that lie beyond the present admissible field.

Second, the framework is explicitly lane-relative in the following sense. Selective locality holds within a constitutive regime, that is, within the space of admissibility fixed by constitutive constraint. The specification of that regime is prior work, discharged by the Constitutive Law. This division of labor is not a concession; it is the correct reflection of the two-role architecture defended in Section 3. The selective role need not explain why reality admits the constitutive regime it does. It needs only to explain, from within that regime, how one admissible continuation rather than another becomes integrated into realized history.

The look-ahead objection therefore fails on inspection. It presupposes that the selective role operates as a dynamics on a pre-existing state space and would require forward simulation to do its work. Neither presupposition survives the distinction between constitutive and selective explanatory roles as drawn in Section 3. What the Localizability Requirement permits is the reading of structural features present in the candidate's position in the admissible order. What it forbids is dependence on features not so present. Intrinsic generative affordance, properly understood, belongs to the first category.

\subsection{Commitment}

Reality must preserve realized history. What has become actual cannot simply vanish from ontological account. This gives the Selective Law a path-dependent context without requiring global foresight: each continuation is judged relative to a cumulative inherited world. Process-oriented metaphysics places exactly this cumulative character of becoming at center stage \citep{Whitehead1929,Rescher1996}.

\begin{definition}[Commitment]
A selective order exhibits commitment when realized history is cumulative: what has become actual remains part of the world's history and conditions future actuality.
\end{definition}

Whitehead matters here for more than the generic thought that reality is processual. His account of concrescence is a nearby analogue of what this paper calls coherent integration, and his contrast between eternal objects and actual occasions parallels, without matching, the present distinction between constitutive admissibility and selective actualization. The difference is equally important. The present paper does not adopt Whitehead's full categoreal scheme and isolates the admissibility question more sharply than Whitehead's own presentation does. Whitehead is therefore engaged as a serious neighboring framework for cumulative becoming rather than as a mere honorary citation.

Commitment does not mean that every local form persists forever. It means that realized actuality is not ontologically reset at each step.

\subsection{Admissibility}

Only continuations coherent with both the Constitutive Law and realized history may be candidates for actualization. The Selective Law therefore acts not on arbitrary possibilities, but on a local continuation cone carved out by coherence.

\begin{definition}[Admissibility]
A continuation of history is admissible when it is coherent with both the Constitutive Law and the realized past it would extend.
\end{definition}

Admissibility is therefore the point at which the constitutive and selective roles touch. Constitutive constraint defines the admissible field; selective actualization acts only within it.

\subsection{Local valuation}

Suppose now that multiple admissible continuations are available. If the Selective Law contains no ordering principle, then actuality is brute. But the relevant ordering cannot be framed as a score assigned by inspecting remote unactualized futures.

A dynamic world must avoid two opposite failures. If it always favors bare preservation, it stagnates. If it always favors novelty as such, it explodes into unstable excess. Hence valuation must be dual. It must regard both what a continuation immediately opens and what it requires in order to become real without destabilizing what already exists. This balance between novelty and stability is also one reason process traditions resist reducing becoming either to mere repetition or to unconstrained flux \citep{Whitehead1929,Rescher1996,Cartwright1999}.

\begin{definition}[Intrinsic generative affordance]
The intrinsic generative affordance of a continuation is the degree to which its present structure immediately opens further coherent continuation space.
\end{definition}

Intrinsic generative affordance is not a forecast produced by simulating remote futures. It is the locally readable generative room already encoded in the continuation itself and in the immediate continuation field it opens.

\begin{definition}[Stabilization burden]
The stabilization burden of a continuation is the additional coordination, integration, and maintenance required to make it actual without destabilizing realized history.
\end{definition}

Intrinsic generative affordance and stabilization burden are both structural properties of a candidate's position in the present admissible field. Neither requires forecasting which admissible continuations will subsequently be integrated. Intrinsic generative affordance without stabilization collapses into uncontrolled branching. Stabilization without intrinsic generative affordance collapses into sterile conservation. A genuine selective rule must therefore weigh both, reading each from features available in realized history, in the candidate itself, and in the immediate structural relations that the candidate's integration would determine.

\subsection{Critical sufficiency}

Even dual valuation is not enough. A world that always chose the abstractly largest continuation would collapse into gratuitous composite leaps. Distinct admissible gains could be bundled into one oversized transition, and actuality would lose articulation.

Therefore the Selective Law cannot be a rule of simple maximization. It must be \emph{critical}: actuality should prefer what is sufficient \emph{here and now}, not what is merely largest in the abstract.

\begin{definition}[Critical sufficiency]
A continuation is sufficient in context when its intrinsic generative affordance is strong enough, relative to its stabilization burden, to justify actualization at the present stage of history.
\end{definition}

Critical sufficiency introduces a context-sensitive threshold. Not every positive continuation actualizes. Only those strong enough to justify integration under present conditions can do so.

Intrinsic generative affordance and stabilization burden are the inputs to valuation. Critical sufficiency is the threshold relation that says when their balance is enough for actualization here and now. It is therefore a scale-condition on eligibility, not a second name for generative affordance.

\subsection{Constructive primeness}

Among sufficient continuations, actuality must prefer those that are not merely artificial bundles of several independently sufficient moves. Otherwise critical sufficiency still leaves room for arbitrary packing, and history again loses articulation.

Constructive primeness answers a different question from critical sufficiency. Critical sufficiency asks whether a candidate is enough. Constructive primeness asks whether the candidate's being enough is genuine or merely the result of packaging together already-enough submoves. The former blocks pressure toward ever larger moves; the latter blocks arbitrary bundling among moves that already meet the threshold.

\begin{definition}[Constructive primeness]
A continuation at a stage fails to be constructively prime when it can be analyzed as a composition of independently sufficient components, each of which would already clear the sufficiency threshold in context on its own. A continuation is constructively prime when it admits no such analysis: its sufficiency is genuine rather than parasitic on the bundling of components that could have been actualized separately.
\end{definition}

Two features of this definition deserve comment.

First, the definition is stated regulatively rather than algebraically. It does not presuppose a specific combining operation on continuations, a specific partial order on decompositions, or a unique-factorization theorem for admissible candidates. It characterizes constructive primeness by what it excludes, gratuitous bundling of independently sufficient components, rather than by a positive mathematical structure. This is deliberate. The foundations paper operates at a level of abstraction that should not legislate the algebra of combination. That algebra is the proper work of implementation.

Second, the regulative constraint is nevertheless determinate enough to do its argumentative work. What constructive primeness blocks is clear: the actualization of composite candidates that clear the sufficiency threshold only because they package together smaller candidates each of which was already sufficient on its own. What it permits is equally clear: the actualization of candidates whose sufficiency is not thus decomposable, whether because the candidate is structurally atomic in the relevant context or because its sufficiency arises from the interaction of components that would not be independently sufficient. The later necessity argument defends the need for this constraint by showing what goes wrong without it.

The mathematical content of ``composition,'' ``component,'' ``independently sufficient,'' and ``analysis'' is left open. Any adequate implementation of the Selective Law must supply these notions within whatever algebraic setting it uses to represent admissible continuations and their internal structure. Different implementations may supply them differently. The foundations paper commits only to the regulative claim that the implementation-level notions, however specified, must respect the exclusion of gratuitous bundling that this definition articulates.

\subsection{Coherent integration}

Actualization must not end with choice alone. What is actualized must enter history in a way that changes the next field of admissible continuations. Without that, there is no cumulative reality, only repeated isolated choices.

\begin{definition}[Coherent integration]
A selective order exhibits coherent integration when the actualized continuation is irreversibly incorporated into realized history and thereby reshapes the future admissible field.
\end{definition}

Commitment and coherent integration are related but not identical. Commitment names the persistence of already realized history as a constraint on later actuality. Coherent integration names the update by which a newly actualized continuation enters that persistent history and reshapes the next admissible field. The former looks backward from the present context; the latter looks forward from the present incorporation.

Coherent integration turns local actualization into world-history rather than episodic choice. Deeper future structure then arises by recursive iteration of the same local law, not by precomputation of a fully surveyed future.

\subsection{Canonical statement of the Selective Law}

We may now state the selective role in canonical form. The statement is deliberately regulative rather than algorithmic: it characterizes what any adequate selective law must achieve, not a specific mathematical operation on a specified algebraic structure.

\begin{theorem}[Selective Law]
Among the admissible continuations of realized history, actuality proceeds by integrating a continuation that is sufficient in context and constructively prime in context, and that is not strictly outranked under the history-indexed preorder by another admissible continuation that is likewise sufficient and prime. The continuation so integrated is irreversibly incorporated into realized history, thereby updating the admissible field for further selection.
\end{theorem}

The statement is regulative in the following sense. It specifies the conditions any adequate selection at a stage must satisfy, admissibility, sufficiency in context, primeness in context, and non-dominance under the history-indexed preorder, and it specifies the integrative effect of selection on realized history. It does not specify, at the foundational level, a unique continuation as the one actual continuation. Among continuations satisfying all four conditions, further selection between them is the work of implementation-level structure that the foundational framework does not itself supply.

\begin{remark}
The phrasing ``proceeds by integrating'' rather than ``selects'' is deliberate. ``Selects'' suggests a single deterministic operation with a unique output, which would require more at the foundational level than the framework supplies. ``Proceeds by integrating'' describes what happens, an admissible continuation satisfying the regulative conditions becomes part of realized history, without overclaiming determinacy. An implementation-level account, such as the one developed in subsequent work under the Principle of Efficient Novelty, may secure determinacy by additional mathematical structure; the foundational framework does not require that it do so.
\end{remark}

This scoping is not a gap in the argument. It reflects the proper division of labor between foundational explanation and formal implementation. The foundational framework characterizes what any adequate selective law must achieve, the structural constraints that render a law intelligible as a law of selective actualization rather than as a brute selector or a globally foresighted oracle. The specification of a unique actual continuation among those satisfying these constraints requires additional structure: a specific algebra of combination on continuations, a specific refinement of the comparative relation, a specific tie-breaking discipline, or a specific valuational metric. These are implementation choices. Different implementations may make different choices and still count as adequate articulations of the same underlying selective role. The underdetermination they introduce is not a third role smuggled in by the back door; it is the familiar difference between a foundational characterization of selective actualization and a complete implementation of it.

\subsection{Why these requirements are needed}

The previous subsections stated the selective role positively. The present subsection strengthens the case by showing that each component is forced by the failure of nearby alternatives.

\subsubsection{Without localizability: teleology or opacity}

If the selective law ordered continuations by properties not readable from the present admissible field, by dependence on continuations beyond that field, or on how far-downstream admissible developments would actually unfold, it would have to proceed either by an implicit teleological formulation or by an opaque global operation not grounded in present structure. Neither is compatible with the explanation of a determinate, historically articulated world by features already present at the stage of actualization. Therefore localizability, in the sense of the Localizability Requirement, is necessary.

\subsubsection{Without commitment: no evolving world}

A world without commitment is not an evolving world but a sequence of disconnected replacements. Nothing that has happened genuinely belongs to what comes next. Identity through change disappears, and with it any meaningful notion of history. Therefore commitment is necessary.

\subsubsection{Without admissibility: self-violation}

A selective law without admissibility can actualize what the constitutive law has already ruled out. In that case the ontology contradicts itself: possibility is constrained at one level and unconstrained at another. Therefore admissibility is necessary.

\subsubsection{Without valuation: brute actuality}

If multiple coherent continuations are available and nothing orders them, then actuality is brute. The world becomes one possibility rather than another for no internal reason. Therefore valuation is necessary.

\subsubsection{Without intrinsic generative affordance: stagnation}

A world that values only preservation may remain stable, but it cannot be genuinely dynamic. It can only conserve what already exists. No principled route to novelty remains. Therefore valuation cannot reduce to stabilization burden alone.

\subsubsection{Without stabilization burden: explosion}

A world that values only immediately opened continuation space may constantly introduce difference, but at the cost of losing stability and identity. Novelty becomes self-undermining. Therefore valuation cannot reduce to intrinsic generative affordance alone.

\subsubsection{Without critical sufficiency: oversized jumps}

If actuality had no context-sensitive threshold, it would have no principled reason to stop at the first continuation that is adequate for actualization. Even a non-composite continuation that is already enough could be bypassed in favor of a larger admissible continuation simply because the latter carries more total affordance. Actualization would then be biased toward delay and overshoot: the world would wait for, or prefer, ever larger moves before integrating anything. History would advance by scale-driven escalation rather than articulated becoming. Therefore the law must be critical rather than naively maximal.

\subsubsection{Without constructive primeness: gratuitous bundling}

Suppose critical sufficiency is retained but constructive primeness is not. Then the world can stop once enough has been reached, but it still cannot distinguish an authentically sufficient move from a package of several already sufficient moves taken together. The failure is no longer one of scale but of composition. Once the threshold is met, arbitrary bundling remains permitted. History would still lose articulation because several independently eligible transitions could be collapsed into one gratuitous aggregate. Therefore constructive primeness is necessary in addition to, rather than instead of, critical sufficiency.

\subsubsection{Without coherent integration: no world-history}

Suppose a continuation is actualized but not coherently integrated. Then it does not alter the future admissible field. Actualization becomes episodic rather than cumulative, and no world-history forms. Therefore coherent integration is necessary.

\subsubsection{Therefore}

Any universe rich enough to sustain determinate, cumulative becoming must obey a Selective Law that is:

\begin{itemize}
    \item cumulative,
    \item coherence-preserving,
    \item locally evaluable,
    \item valuational,
    \item critical,
    \item prime-selective,
    \item and integrative.
\end{itemize}

Taken together, these constraints force selective actualization into something close to the canonical statement given above. The exact later formalization may vary. The structural shape cannot vary far without the ontology collapsing into arbitrariness, stagnation, explosion, serial replacement, or computational opacity.

\section{A Minimal Formal Schema}

The preceding argument can now be compressed into a minimal formal schema. The purpose of the schema is not to replace the philosophical discussion with a technical system. It is to state, in compact form, what the selective role ranges over, how context enters, what counts as a better continuation, and how an actual choice becomes the next realized history.

The notation is intentionally spare:

\begin{itemize}
    \item $H$ is realized history up to the present stage.
    \item $A(H)$ is the admissible continuation cone at $H$: the continuations compatible with constitutive constraint and with the commitments already carried by $H$.
    \item $\preceq_H$ is a context-sensitive preorder on $A(H)$ used here as a convenient schematic representation of local comparative eligibility relative to $H$, once the local balance of generative affordance and stabilization burden is taken into account.
    \item $S_H(x)$ says that $x$ is sufficient at $H$: it clears the context-sensitive threshold for actualization.
    \item $P_H(x)$ says that $x$ is prime at $H$: it is not merely a gratuitous bundle of independently sufficient continuations.
    \item $I(H,x)$ is the integration of $x$ into the next realized history.
\end{itemize}

The dependence on $H$ is not cosmetic. It represents the path-dependent context inherited from realized history. Admissibility, sufficiency, primeness, and comparative rank are all assessed relative to the history being extended. The framework therefore does not compare continuations against a history-free global scale. It compares admissible continuations relative to a determinate present continuation cone.

The comparative structure is kept abstract on purpose. The philosophical argument up to this point forces some local comparative articulation, but not a preorder in particular. A preorder is used here because it is a familiar and minimally committal device for representing non-dominance, context-sensitivity, and ties without presupposing a metric or a total ranking. A later implementation could instead employ a partial order, a choice correspondence over an eligible set, a partition into sufficient and insufficient regions supplemented by local dominance rules, a similarity structure, or a non-additive ordering. What the foundations paper requires is only enough local comparative structure to distinguish dominated from non-dominated admissible continuations in context $H$.

Using a preorder rather than a total order also leaves room for ties. Two admissible continuations may be equally eligible at this level of abstraction. A later theory may refine such ties by additional local structure, treat them as genuine parity, or assign them deterministic, probabilistic, or amplitude-like refinements. Such refinements do not add a third foundational role. They are more determinate implementations of the selective role. What matters here is that ``better'' means better relative to the realized history and its admissible continuation cone, not better by inspection of indefinitely distant futures.

The preorder $\preceq_H$ ranks admissible continuations by structural features readable at the present stage: features of the candidate itself, of the realized history it would extend, and of the immediate admissible field its integration would determine. It does not rank by forecasting which admissible continuations will be integrated at subsequent stages.

With that notation in place, the selective role can be stated schematically. A continuation $x$ is eligible for actualization at $H$ only if it lies in the admissible cone, is sufficient in context, is prime in context, and is not strictly outranked by another admissible continuation that is likewise sufficient and prime:

\[
\begin{aligned}
&x \in A(H), \qquad S_H(x), \qquad P_H(x),\\
&\neg \exists y \in A(H)\bigl(S_H(y) \wedge P_H(y) \wedge y \prec_H x\bigr).
\end{aligned}
\]

Here $y \prec_H x$ is the strict part of the preorder: $y$ is strictly more eligible than $x$ in context $H$. Once some such $x$ is integrated, the next realized history is given schematically by

\[
H' = I(H,x).
\]

This schema is minimal but already clarifying. It says what an adequate implementation must integrate: an admissible continuation satisfying sufficiency, primeness, and non-dominance in context. It says what counts as better: comparative eligibility under the history-indexed preorder $\preceq_H$. It says how context enters: every condition is indexed to the realized history $H$ being extended. And it says how actuality advances: not by cataloguing possible worlds, but by integrating one locally admissible, sufficient, and prime continuation into the next history. That is enough structure to make the framework assessable without importing a full technical implementation.

\section{A Toy Branching Model}

A small branching example makes the division of explanatory labor fully explicit. Let $H$ be realized history, and consider four candidate continuations:

\[
a, \qquad b, \qquad c, \qquad d = b \oplus c.
\]

Here $\oplus$ is schematic notation for whatever combining structure an implementation supplies when representing how one continuation can be analyzed as incorporating others. This paper does not commit to a specific algebraic reading of $\oplus$, because the algebra of combination is proper implementation work. What the example requires is only that some such structure be available when the framework is formally implemented, so that the regulative notion of constructive primeness can be given determinate content. The toy model illustrates the architecture of the view at the regulative level, and the schematic character of $\oplus$ reflects the presentation-independence of that level rather than an unresolved commitment.

The case of $a$ illustrates constitutive exclusion. Suppose $a$ carries an unresolved incompatibility: it cannot be jointly determined with the identity and coherence conditions already borne by $H$. Then $a$ never enters the admissible continuation cone at all:

\[
a \notin A(H).
\]

This is the work of constitutive constraint. The selective role does not need to reject $a$, because $a$ was never a genuine candidate for actualization in the first place.

The case of $b$ illustrates admissibility without sufficiency. Suppose $b$ is coherent with constitutive constraint and with realized history, so that

\[
b \in A(H),
\]

but suppose also that $b$ opens too little local generative affordance to justify the stabilization burden it would impose in the present context. Then $b$ is admissible but not sufficient:

\[
b \in A(H), \qquad \neg S_H(b).
\]

This is the first point at which the selective role does substantive work. Constitutive constraint leaves $b$ inside the admissible cone, but selective actualization does not yet permit its integration into realized history.

The case of $c$ illustrates a continuation that is admissible, sufficient, and prime. Suppose $c$ is coherent with $H$, opens enough further continuation space to justify its stabilization burden, and does so without being a gratuitous bundle of independently sufficient parts. Then

\[
c \in A(H), \qquad S_H(c), \qquad P_H(c).
\]

To keep the example focused, suppose further that no admissible continuation that is both sufficient and prime strictly outranks $c$ under the context-sensitive preorder:

\[
\neg \exists y \in A(H)\bigl(S_H(y) \wedge P_H(y) \wedge y \prec_H c\bigr).
\]

Under those assumptions, $c$ is eligible for actualization.

The final case shows why primeness matters. Let

\[
d = b \oplus c.
\]

Assume that this bundled continuation is still admissible and still sufficient because it contains the gain already carried by $c$:

\[
d \in A(H), \qquad S_H(d).
\]

Even so, $d$ should not be selected if the extra $b$-component contributes nothing that is contextually needed for actualization. In that case $d$ is not prime. It is merely a gratuitous package that combines one independently insufficient move with one move that was already sufficient on its own:

\[
\neg P_H(d).
\]

The example therefore yields a clean pattern of exclusion and selection:

\[
\begin{aligned}
&a \notin A(H),\\
&b \in A(H)\ \text{but}\ \neg S_H(b),\\
&c \in A(H) \wedge S_H(c) \wedge P_H(c),\\
&d \in A(H) \wedge S_H(d) \wedge \neg P_H(d).
\end{aligned}
\]

Only $c$ survives every filter that the framework imposes. Constitutive constraint excludes $a$. Selective sufficiency excludes $b$. Primeness excludes $d$. Under the assumptions made explicit in each case, one may therefore write schematically

\[
H' = I(H,c).
\]

The toy model displays the architecture of the view in one place. Constitutive constraint determines which continuations count as genuine candidates. Selective actualization then distinguishes the merely admissible from the actually eligible via sufficiency, and distinguishes the bundled from the constructively prime via the exclusion of gratuitous packaging. The result, under the assumptions made explicit in each case, is that $c$ is the continuation that satisfies all the regulative conditions of Theorem~2. That $c$ rather than any other equally qualifying candidate becomes integrated into realized history, in settings where multiple such candidates would coexist, is a question the foundational framework does not on its own settle. Its settlement belongs to the implementation-level structure that a specific articulation of the Selective Law supplies. The architectural point of the toy model is not that $c$ is uniquely determined by the framework; it is that the framework determines what must be true of whatever continuation is integrated, and that this is enough to exhibit the two-role division of explanatory labor.

\section{What the Framework Adds}

\subsection{A grounding-theoretic translation}

A further way to locate the contribution is by translating it into the vocabulary of the contemporary grounding literature without thereby committing the paper to any particular position within it. On that vocabulary, the constitutive role corresponds to a relation of structural grounding: coherence, equivalence-invariance, canonicity, and local satisfiability of compatibility burdens jointly ground the facts that determine admissibility. The selective role corresponds to a relation of partial grounding on admissibility facts: given that a continuation is admissible, its integration into realized history is grounded in its satisfaction of sufficiency, primeness, and non-dominance in context, but only partially, since among continuations satisfying all the regulative conditions the foundational framework does not uniquely ground which is integrated. That further determination belongs to implementation-level structure, and is therefore grounded at a level the foundational framework defers rather than adjudicates \citep{Fine2012,Schaffer2009,Rosen2010}.

The three-category classification defended in subsection 3.3 admits a similar translation. Roles that are structural consequences of how the two foundational roles operate, modal structure being the paradigm case, stand in a relation of derivative grounding: the grounds of modal facts are the joint grounds of admissibility facts and selective facts, so modal facts are grounded at a lower level than either constitutive or selective facts taken alone. Roles that are derived features of cumulative actualization, individuation being the paradigm case, stand in a relation of iterated path-dependent grounding: individuation facts are grounded in the cumulative history of applications of the selective role, each of which partially grounds the subsequent stage. Scope questions about what grounds the framework itself are ungrounded in the relevant sense: they do not admit of further grounding without presupposing the framework whose grounding they purport to specify, and their status is analogous to that of fundamental primitives in standard grounding-theoretic treatments.

This translation is offered as a bridge rather than as an endorsement of any particular grounding theorist's view. A Finean account of grounding, a Schafferian account, and a Rosen-style account each admit the bridge, with differences in how they articulate the relevant grounding relations. What matters for the present paper is that the architectural thesis, two irreducible foundational roles with a principled classification of apparent third roles, is stateable in grounding vocabulary and is therefore assessable by grounding-literate readers without requiring that the paper adopt a particular grounding theory to make its case.

\subsection{Order of form and order of becoming}

Constitutive constraint and selective actualization are not rival explanations. They are complementary.

Constitutive constraint answers the question: \emph{what can coherently be?} It defines the admissible order of form. In a broadly structural register, it concerns the realm of coherent candidates as such.

Selective actualization answers the question: \emph{what becomes actual?} It defines the order of realized history through recursive local actualization. In a broadly process-oriented register, it concerns the transition from candidate structure to inhabited actuality.

A useful way to state the contribution is this: constitutive constraint gives the world its \emph{order of form}, while selective actualization gives it its \emph{order of becoming}. The novelty is not merely to contrast possibility with actuality in the abstract, but to argue that foundational explanation must assign these two questions to distinct explanatory roles. The first specifies the candidate space but does not order actualization within it. The second orders actualization within the candidate space but does not create that space. Neither is reducible to the other without explanatory loss. In that respect the framework remains compatible with sharply different views of lawhood and structure while pressing a more basic question about the relation between admissibility and actuality \citep{Armstrong1983,vanFraassen1989,Maudlin2007,French2014}.

\section{Scope, Limits, and What Is Deferred}

A foundations paper gains strength by distinguishing necessity from proposal and proposal from speculation.

\subsection{Argued as metaphysical necessity}

The following claims are argued here as necessary for any stable and historically articulated world:

\begin{enumerate}
    \item Reality requires both a constitutive law and a selective law.
    \item The constitutive law must at least require coherent realizability, equivalence-invariance, canonical executability, and local satisfiability of compatibility burdens.
    \item The selective law must at least require localizability, commitment, admissibility, valuation, critical sufficiency, constructive primeness understood regulatively, and coherent integration.
\end{enumerate}

These are the core philosophical results of the paper.

\subsection{Proposed as canonical formulation}

The following stronger claims are proposed as the best current articulation of the foregoing necessities:

\begin{enumerate}
    \item The Constitutive Law is best stated as the requirement that a universe of the target kind be a coherent, equivalence-invariant, canonically executable structure whose identity and compatibility conditions are locally satisfiable.
    \item The Selective Law is best stated as the regulative rule that actuality proceeds by integrating a constructively prime admissible continuation whose intrinsic generative affordance is sufficient to justify its stabilization burden in context, and that the continuation so integrated irreversibly updates the basis of further actualization. The rule characterizes what any adequate implementation must achieve rather than uniquely specifying which continuation is actualized.
\end{enumerate}

These formulations are stronger than the bare necessities because they package the minimal requirements into single canonical statements.

\subsection{Deferred to later work}

The following questions are intentionally left open here:

\begin{enumerate}
    \item The exact mathematical form of valuation and local actualization.
    \item Any concrete evaluator, threshold equation, or dynamical mechanism that might instantiate the Selective Law.
    \item The formal status of arithmetic and the natural numbers, beyond the weak assumption that cumulative actualization can be discussed stagewise.
    \item Probabilistic or amplitude-like refinements of actualization.
    \item Cosmological, physical, or phenomenological applications.
    \item Any claim that a single human formalism uniquely and exhaustively captures the constitutive order.
    \item The specification of implementation-level structure, including the algebra of combination on continuations, the choice of local comparative structure represented schematically here by $\preceq_H$, and the deterministic, probabilistic, or otherwise controlled discipline by which uniqueness or non-uniqueness of actualization is handled among continuations satisfying the regulative conditions of Theorem~2, is deferred to formal work on specific articulations of the Selective Law, of which the Principle of Efficient Novelty is one example.
\end{enumerate}

This paper is therefore not a complete metaphysical system. It is the identification of the two irreducible tasks any such system must perform, together with the strongest presentation-independent articulation currently available.

\section{Conclusion}

The paper has argued for a simple but demanding thesis. A universe cannot be constituted by form alone, because form alone never explains actuality. Nor can it be constituted by actuality alone, because actuality without admissible form is unintelligible. A determinate world requires both a constitutive law and a selective law.

The constitutive order determines what can coherently count as possible. The selective order determines how coherent possibilities are locally actualized and integrated into history. At the indexed level of foundational explanation defended in subsection 3.3, no third irreducible role has yet survived examination; the classification of modal structure, individuation, and scope questions given there is offered as a structural proposal open to further challenge. Every further putatively primitive law considered so far has either constrained coherent possibility, governed actualization from within it, or fallen into one of the derivative or scope-sensitive categories distinguished there.

The resulting picture may be summarized in one sentence:

\begin{quote}
Reality is neither mere form nor brute fact. It is coherent possibility under a law of local actualization.
\end{quote}

Or, still more compactly:

\begin{quote}
A universe of the target kind exists only where coherent form and local actualization meet.
\end{quote}

That is the paper's final claim. If it is right, then the deepest problem for metaphysics is not to choose between mathematics and world, nor between possibility and actuality, but to understand why reality requires both -- and why the two laws that govern them must have the form they do.

\bibliographystyle{plainnat}
\bibliography{foundations_reality_refs}

\end{document}
